\clearpage
\SourceLeaf{398}
\fancyhead[L]{\small SUPPLEMENTS TO THE GEOMETRY OF POSITION.}
\pdfbookmark[1]{Supplements to Carnot's Geometry of Position. 1810.}{carnot-zusaetze}
\begin{center}
{\Large SUPPLEMENTS.}\par\bigskip
TO CARNOT'S GEOMETRY OF POSITION\par\bigskip
TRANSLATED BY SCHUMACHER.\quad 1810.\par\bigskip
\rule{.1\textwidth}{.3pt}
\end{center}
\par\medskip\Needspace{5\baselineskip}
\pdfbookmark[2]{I.}{carnot-zusatz-i}
\centerline{\textbf{I.}}\nobreak\smallskip
\{I owe the following analytical treatment of the notable points of a triangle to the kindness of Professor \textsc{Gauss}. \textsc{Schumacher}.\}

Let \(A,A',A''\) be the three vertices of a triangle and let their respective coordinates be
\[
\begin{array}{ll}
x,&y\\x',&y'\\x'',&y''
\end{array}
\]
\par
The coordinates of \(B,B',B''\) that bisect the sides \(A'A'',A''A,AA'\) are evidently
\[
\begin{array}{ll}
\tfrac12(x'+x''),&\tfrac12(y'+y'')\\
\tfrac12(x''+x),&\tfrac12(y''+y)\\
\tfrac12(x+x'),&\tfrac12(y+y')
\end{array}
\]
\par
On the lines \(AB,A'B',A''B''\), extending them forward or backward if necessary, take segments measured from \(A,A',A''\) proportional to the corresponding vertex-to-midpoint segments, in the ratio \(n:1\). If the segments are taken backward, \(n\) is to be regarded as negative. Let \(C,C',C''\) be the endpoints of these segments. Their coordinates are
